\section{Problem Setup and Protocol}
\label{sec:setup}

\subsection{Speckle forward model and information support}
We simulate the memory-effect regime in which a thin diffuser converts coherent illumination into a shift-invariant speckle intensity point-spread function (PSF). Working in one dimension of notation, the measurement is
\begin{equation}
y = h \ast x + n, \qquad n \sim \text{Poisson--Gaussian},
\label{eq:forward}
\end{equation}
where $x$ is the object, $h$ is the intensity PSF generated by a circular pupil of radius $r$ with a random phase screen, and $\ast$ denotes convolution. The system's \emph{information support} is the half-power radius $R_{\MTF}$ of the modulation transfer function (MTF): spatial frequencies beyond $R_{\MTF}$ are not transported by the measurement, and any reconstruction energy there is fabricated by the model. For the pupil model used throughout, $R_{\MTF} = 2r$ pixels analytically; we verified empirically that single-realization MTF half-power estimation is unusable (speckle-dominated), which is why the protocol uses the analytic support and the object-bandwidth factors are expressed in units of $\sigma_{\MTF} = R_{\MTF}/2.355$.

\subsection{Endpoints}
\paragraph{Out-of-band energy (fabrication proxy).}
For a reconstruction $\hat{x}$,
\begin{equation}
\oob(\hat{x}) = \frac{\sum_{|\nu| > R_{\MTF}} |\hat{X}(\nu)|^2}{\sum_{\nu} |\hat{X}(\nu)|^2},
\label{eq:oob}
\end{equation}
the fraction of reconstruction spectral energy outside the measured information support. It quantifies fabrication beyond what the measurement can contain. It is a proxy: it does not measure perceptual or semantic error, and a low value is not by itself evidence of a good reconstruction (a constant image scores zero). It is used here as the dependent variable of a controlled intervention, anchored to a measured physical quantity.

\paragraph{Relative forward consistency (no-reference endpoint).}
\begin{equation}
\relfwd(\hat{x}) = \frac{\lVert h \ast \hat{x} - y \rVert_2}{\lVert y \rVert_2},
\label{eq:relfwd}
\end{equation}
the degree to which a reconstruction re-explains its own measurement. It requires no ground truth and is the primary endpoint for checking that fabrication suppression does not come at the cost of measurement agreement.

\paragraph{Band-limited PSNR (secondary).}
PSNR against a ground truth low-pass filtered at the target bandwidth. We report it as a secondary, reference-dependent quantity; Sec.~\ref{sec:baselines} shows it is not innocent: the ranking of models under this metric depends on the bandwidth of the reference used, which is one reason the protocol treats it as secondary.

\subsection{The bandwidth intervention}
All arms share the measurement set, the architecture, the initialization, and the optimizer. The intervention changes only the bandwidth of the training (or fine-tuning) target:
\begin{itemize}
\item \textbf{raw target} ($\sigma = 0$): the unfiltered ground-truth object;
\item \textbf{matched target} ($\sigma = 0.9\,\sigma_{\MTF}$): ground truth low-pass filtered to the system's measured support;
\item \textbf{half target} ($\sigma = 0.45\,\sigma_{\MTF}$): a conservatively band-limited target.
\end{itemize}
The central comparison is a \emph{paired} one: from a converged raw-target baseline (20 epochs), one arm continues on the same raw target for four additional epochs (the \textbf{compute-matched control}), while the other fine-tunes on a band-limited target for the same four epochs. The control isolates the bandwidth change from the effect of extra optimization, which is the primary confound of any fine-tuning comparison.

\subsection{Architectures, budgets, and decision rules}
The decoder is a three-scale residual U-Net with a global skip connection, trained with $L_2$ loss, Adam, gradient clipping, and identical seeds across arms; two widths (base 16 and base 32) are tested with three seeds each, on $64\times64$ images with 2{,}400 training and 400 test objects (sharpened rendered digits). The residual skip removes the all-zero output attractor observed in early pilots and is used in all reported results. The main fine-tune comparison (Sec.~\ref{sec:main-control}) was pre-registered with the decision rule: causal effect if the fine-tuned-to-control out-of-band ratio is below $0.8$; the suppression direction is additionally tested by an exact one-sided sign test over paired units. Quantities are reported as mean $\pm$ standard deviation over the six units (2 architectures $\times$ 3 seeds) unless stated otherwise. All results in Secs.~\ref{sec:experiments}--\ref{sec:svpsf} come from this simulated protocol; Sec.~\ref{sec:realdata} treats real data separately.
